\documentclass[11pt]{article}

\usepackage[margin=1in]{geometry}
\usepackage{amsmath,amssymb,amsthm}
\usepackage{graphicx}
\usepackage[dvipsnames,table]{xcolor}
\usepackage[round,authoryear]{natbib}
\usepackage{hyperref}
\usepackage{booktabs}

\hypersetup{colorlinks=true,linkcolor=MidnightBlue,citecolor=MidnightBlue,urlcolor=NavyBlue}

\graphicspath{{figures/}}

\newtheorem{theorem}{Theorem}
\newtheorem{corollary}{Corollary}
\newtheorem{lemma}{Lemma}

\newcommand{\vth}{v_{\mathrm{th}}}
\newcommand{\gradpar}{\nabla_\parallel}
\newcommand{\kperp}{k_\perp}
\newcommand{\zpm}{z^{\pm}}

\title{Kinetic reduced magnetohydrodynamics as a continuum of passive scalars in phase velocity}
\author{A. Kanekar}
\date{Draft, \today}

\begin{document}
\maketitle

\begin{abstract}
% Written last.
\end{abstract}

\section{Introduction}
\label{sec:intro}
% Compressive fluctuations in KRMHD; phase mixing, nonlinear advection and echoes;
% the invariants of Chandran, Mallet & Meyrand (2026); what this paper shows.

\section{The compressive sector of KRMHD}
\label{sec:model}
% CMM Appendix B form: two branches G^\pm with \Lambda^\pm, the Laguerre l >= 2 sector,
% fluid moments, conventions (z^+ direction, v_th, s = \sqrt{\beta_i}).

\section{Phase-velocity slices}
\label{sec:slices}
% The operator L_\Lambda, the free-energy inner product, its spectrum, the Case--Van Kampen
% transform, and the slice equation: each slice is a passive scalar advected by
% \Phi + s u \Psi and streamed at s u v_A.

\section{Invariants}
\label{sec:invariants}
% Q_\phi = [G, \phi(L) G]; W, \Gamma, H as \phi = 1, u, 1/u; \Gamma as the field-aligned
% free-energy flux; per-slice invariants; what survives at finite k_\perp \rho_i.

\section{Lagrangian linearity}
\label{sec:frozen}
% The frozen-line theorem and its corollaries for single-point statistics; forcing;
% what breaks it; what it implies for echoes.

\section{Phase mixing as decorrelation of slices}
\label{sec:decorrelation}
% Two-slice law; Hermite-spectrum dictionary and the m^{-1/2} prefactor; the Hermite flux;
% the sgn(k_z) decomposition under field-line tilt.

\section{Symmetry and special slices}
\label{sec:symmetry}
% Mirror symmetry and balanced turbulence; the s u = \pm 1 slices; u = 0; the fluid limit.

\section{Truncated Hermite hierarchies}
\label{sec:truncated}
% M+1 slices; eigenvalues and u_max(M); CFL; the zero slice at even M; hypercollisions as
% the only coupling between slices.

\section{Checks with GANDALF}
\label{sec:numerics}

\section{Discussion}
\label{sec:discussion}

\bibliographystyle{plainnat}
\bibliography{refs}

\end{document}
